%==============================================================================
% Sjabloon poster bachproef
%==============================================================================
% Gebaseerd op document class `a0poster' door Gerlinde Kettl en Matthias Weiser
% Aangepast voor gebruik aan HOGENT door Jens Buysse en Bert Van Vreckem

\documentclass[a0,portrait]{hogent-poster}

% Info over de opleiding
\course{Bachelorproef}
\studyprogramme{toegepaste informatica}
\academicyear{2025-2026}
\institution{Hogeschool Gent, Valentin Vaerwyckweg 1, 9000 Gent}

% Info over de bachelorproef
\title{Beveiliging van AI-agents in Fintech: automatische detectie van beveiligingslekken in integraties op basis van het Model Context Protocol (MCP)}
\author{Arthur Van Oudenhove}
\email{arthur.vanoudenhove@student.hogent.be}
\supervisor{Lieven Smits}
\cosupervisor{Andres Algaba (Vrije Universiteit Brussel)}

% Indien ingevuld, wordt deze informatie toegevoegd aan het einde van de
% abstract. Zet in commentaar als je dit niet wilt.
\specialisation{AI \& Data Science}
\keywords{MCP, AI, Security, Fintech}
\projectrepo{https://github.com/Arthur-VO/bachelorproef}

%% Extra pakketten voor visualisaties en bronverwijzingen
\usetikzlibrary{positioning}
\usepackage{pgfplots}
\pgfplotsset{compat=1.18}

%% Bibliografie: zelfde stijl en bron als de bachelorproef
\usepackage[style=apa,backend=biber]{biblatex}
\addbibresource{../bachproef/bachproef.bib}
\renewcommand*{\bibfont}{\footnotesize}

\begin{document}

\maketitle

\begin{abstract}
Autonome AI-agents verwerken in de Fintech-sector steeds vaker betalingstransacties en andere bedrijfskritische taken. De communicatie met interne systemen verloopt via het Model Context Protocol (MCP), een standaard die integratie vergemakkelijkt maar tegelijk nieuwe aanvalsvlakken opent: \emph{prompt injection}, \emph{tool poisoning} en \emph{privilege escalation}. Generieke tools zoals Burp Suite en Semgrep missen de contextuele kennis om MCP-specifieke kwetsbaarheden systematisch op te sporen, en handmatige code reviews zijn arbeidsintensief en foutgevoelig. Dit onderzoek ontwikkelt een geautomatiseerd security assessment framework en vergelijkt het met een handmatige beoordeling op een opzettelijk kwetsbare testserver (\emph{Goat}). Het framework detecteerde alle vijf ingebedde kwetsbaarheden in minder dan 30~seconden bij een precisie van minstens 92,3\%; de handmatige procedure vond er vier in 47~minuten.
\end{abstract}

\begin{multicols}{2} % This is how many columns your poster will be broken into, a portrait poster is generally split into 2 columns

\section{Introductie}

Autonome AI-agents nemen in de Fintech-sector steeds meer bedrijfskritische taken over: ze categoriseren rekeningtransacties, initiëren betalingen en krijgen daarvoor directe toegang tot interne systemen \autocite{russell2021artificialintelligence}. Het Model Context Protocol (MCP) is in korte tijd breed geadopteerd als gestandaardiseerde interface tussen AI-modellen en bedrijfsapplicaties \autocite{anthropic2024mcpwhitepaper}.

Die standaardisatie heeft een keerzijde op het vlak van beveiliging. Aanvallers kunnen via \emph{prompt injection} de instructies van de agent manipuleren, via \emph{tool poisoning} schadelijke commando's verbergen in datastromen, en via \emph{privilege escalation} onbedoeld verhoogde toegangsrechten verkrijgen \autocite{song2025protocolunveilingattackvectors}.

Generieke beveiligingstools zoals Burp Suite en Semgrep missen de contextuele kennis om deze MCP-specifieke kwetsbaarheden systematisch te detecteren, en handmatige code reviews zijn arbeidsintensief en foutgevoelig. Dat 56\% van de ondernemingen die AI-agents inzetten beveiligingsincidenten rapporteert in hun MCP-omgevingen, maakt de urgentie concreet \autocite{guo2025systematicanalysismcpsecurity}.

\section{Onderzoeksvraag \& aanpak}

\emph{Hoe kan een framework worden ontwikkeld dat MCP-server implementaties automatisch test op een vooraf gedefinieerde set kritieke kwetsbaarheden, en hoe verhoudt de dekkingsgraad en benodigde tijd per assessment zich tot een handmatige beoordeling op basis van een concrete evaluatiecase?}

Het framework combineert statische analyse (SAST) met dynamische fuzzing en draait containergebaseerd via Docker~Compose. Als evaluatiecase dient een opzettelijk kwetsbare testserver (\emph{Goat}) met vijf bekende, ingebedde kwetsbaarheden. Dezelfde server werd parallel handmatig beoordeeld via code review en HTTP-tests als benchmark.

\begin{center}
  \captionsetup{type=figure}
  \begin{tikzpicture}[
      box/.style={draw, rounded corners, very thick, minimum width=9cm,
                   minimum height=4.5cm, align=center, font=\large},
      arr/.style={->, very thick, line width=2pt}
    ]
    \node[box, fill=hogent-pink!30]      (goat) at (0,0)  {MCP\\Goat-server\\\footnotesize (5 kwetsbaarheden)};
    \node[box, fill=hogent-darkgreen!25] (scan) at (11,0) {Scanner\\SAST + fuzzing};
    \node[box, fill=hogent-ochre!35]     (rep)  at (22,0) {Rapport\\\footnotesize (bevindingen)};
    \draw[arr] (goat) -- (scan);
    \draw[arr] (scan) -- (rep);
  \end{tikzpicture}
  \captionof{figure}{De gecontaineriseerde pijplijn: de scanner test de Goat-server statisch en dynamisch en levert een rapport met bevindingen.}
\end{center}

\section{Resultaten}

Het framework detecteerde alle vijf ingebedde kwetsbaarheden (recall 100\%) in minder dan 30~seconden. Van de 13 bevindingen is er één debateerbaar als false positive, wat een precisie van minstens 92,3\% oplevert. De handmatige procedure vond vier van de vijf kwetsbaarheden in 47~minuten en miste de tool poisoning, die enkel bij runtime zichtbaar is.

\begin{center}
\begin{minipage}{\linewidth}
  \centering
  \captionsetup{type=figure}
  \begin{tikzpicture}
    \begin{axis}[
        ybar,
        bar width=1.4cm,
        width=0.95\linewidth,
        height=13cm,
        ymin=0, ymax=115,
        ylabel={\%},
        symbolic x coords={Recall, Precisie},
        xtick=data,
        nodes near coords,
        every node near coord/.append style={font=\footnotesize},
        enlarge x limits=0.5,
        legend style={at={(0.5,-0.10)}, anchor=north, legend columns=-1},
      ]
      \addplot[fill=hogent-blue]  coordinates {(Recall,100)  (Precisie,92.3)};
      \addplot[fill=hogent-ochre] coordinates {(Recall,80)   (Precisie,100)};
      \legend{Framework, Handmatig}
    \end{axis}
  \end{tikzpicture}
  \captionof{figure}{Detectiekwaliteit: het framework evenaart de handmatige beoordeling op recall en precisie.}
\end{minipage}
\end{center}

\begin{center}
  \captionsetup{type=table}
  \begin{tabular}{@{}lll@{}}
    \toprule
    \textbf{Meetdimensie}   & \textbf{Framework} & \textbf{Handmatig} \\ \midrule
    Benodigde tijd          & $<$30 seconden     & 47 minuten \\
    Reproduceerbaarheid     & Gegarandeerd       & Reviewerafhankelijk \\ \bottomrule
  \end{tabular}
  \captionof{table}{Tijd en reproduceerbaarheid: het structurele voordeel van automatisering.}
\end{center}

De scanner is deterministisch: dezelfde omgeving levert altijd dezelfde bevindingen op, ongeacht de reviewer. De fuzzer detecteert bovendien runtime-gedrag dat onzichtbaar is in broncode, zoals de tool poisoning in \texttt{fetch\_market\_news}.

\section{Conclusies}

Een geautomatiseerd assessment framework detecteert kritieke kwetsbaarheden in MCP-serverimplementaties effectief: sneller, reproduceerbaar en met een dekkingsgraad die in deze evaluatiecase de handmatige beoordeling evenaart. Voor de interpretatie van bedrijfsimpact en complexe logische beveiligingsproblemen blijft menselijke analyse onmisbaar.

Aanbevelingen:

\begin{itemize}
  \item Verplicht autorisatiechecks op serverniveau voor elke tool die toegang biedt tot transacties of gebruikersgegevens.
  \item Beschouw tool outputs als onvertrouwde invoer en saniteer ze op dezelfde manier als inputs.
  \item Gebruik het framework als screeningslaag, niet als vervanging van een security audit.
  \item Integreer de scanner in het CI/CD-proces voor continue compliancemonitoring.
\end{itemize}

\section{Toekomstig onderzoek}

De resultaten steunen op één gecontroleerde evaluatiecase, waardoor de externe validiteit beperkt is. Vervolgonderzoek kan het framework evalueren op meerdere, diverse MCP-servers, bijvoorbeeld via de MCPTox-benchmark van 45 echte servers en 353 tools \autocite{wang2025mcptox}. Een tweede spoor is volledige integratie in een CI/CD-pijplijn met automatische rapportage aan het development team. Een derde open vraag is hoe het framework zich gedraagt bij semantische kwetsbaarheden die afhangen van de systeemprompt van het onderliggende taalmodel.

\section{Bronnen}
\printbibliography[heading=none]

\end{multicols}
\end{document}
